\section{问题重述}

\subsection{问题背景}

脑电图（Electroencephalography，EEG）能够以非侵入方式记录大量神经元群体活动所形成的头皮电位变化，是连接基础神经科学、脑机接口和神经精神疾病研究的重要观测手段。与单纯依靠统计学习完成“输入—标签”映射不同，本题更关注脑电信号背后的\textbf{生成机制与认知机制}：一方面，需要从含有强伪影的原始EEG中恢复与视觉刺激有关的微弱响应；另一方面，还需要建立从视觉通路、皮层神经群体活动到头皮电极观测的计算联系，并进一步讨论视觉信息与记忆、认知状态之间的动态耦合。

题目给出两个视觉认知实验项目。实验项目I中，受试者在目标出现前已经知道目标位置，主要等待目标出现；实验项目II中，受试者事先不知道目标位于左侧还是右侧，只知道视觉提示的三角形方向。两个项目均在额区F3、Fz、F4记录EEG，因此同一视觉刺激在不同任务条件下的脑电差异，既可能包含直接视觉响应，也可能混合注意、预期、记忆匹配等高层认知因素。

\subsection{实验数据说明}

题目提供四组实验数据：
\[
\texttt{VisualCogA\_Task-1.mat},\quad
\texttt{VisualCogA\_Task-2.mat},\quad
\texttt{VisualCogB\_Task-1.mat},\quad
\texttt{VisualCogB\_Task-2.mat},
\]
采样频率均为$256\,\mathrm{Hz}$。每组数据包含10个记录通道，其中Fz、F3、F4为原始EEG，视觉提示通道用于标记左、右三角提示，目标应答通道记录实际点击行为，另有时间戳等辅助信息。题目同时给出了设备自带滤波后的FzDecon、F3Decon、F4Decon，但明确指出该处理可能在去除干扰的同时破坏视觉响应特征，因此本文不将其作为主要建模输入，而从原始Fz、F3、F4信号出发独立设计预处理流程。

\subsection{问题一：脑电预处理与有效视觉响应提取}

视觉刺激可能在约$250$--$500\,\mathrm{ms}$形成P300相关晚期响应，但原始EEG同时受到眨眼、运动、慢漂移及高频噪声等多种干扰。问题一要求分别针对实验项目I和项目II设计特定降噪方法，在抑制伪影的同时尽量保留视物形状相关特征，随后提取F3、Fz、F4各观测点的有效视觉响应并拟合相应曲线。

因此，问题一的关键并非“滤得越平滑越好”，而是建立可量化的\textbf{降噪—保真折中准则}，并证明最终处理方法没有因为去噪而系统性削弱潜在视觉成分。

\subsection{问题二：LGN—皮层—头皮EEG生成机制}

问题二要求建立从外侧膝状体核（LGN）到皮层、再到头皮EEG观测的计算模型，并利用该模型解释左、右三角形视觉刺激产生不同脑电信号的形成机制，最终构造能够区分左右视觉刺激响应的特征表示。

该问题本质上要求把“刺激图形的空间结构”与“神经群体响应分布”和“头皮电极观测”连接起来。由于三通道EEG不足以唯一反演真实三维神经源，本文将采用受神经生理约束的低维生成模型，在可辨识性与生理解释之间取得平衡。

\subsection{问题三：视觉认知宏观模型}

问题三要求截取与视觉认知相关联的脑电信号，其终点可位于行为应答前约$100\,\mathrm{ms}$，在问题二脑电形成机制的基础上，同时考虑信号可能来源于海马体、LGN等结构，建立视觉认知的宏观模型，并使用题目提供的EEG进行验证。题目特别指出，应答并不一定正确，也可能出现未及时应答，因此不能直接用行为结果替代认知状态本身。

由此，问题三需要在问题二“直接视觉响应”的基础上进一步引入记忆匹配和认知状态，并把视觉提示、内部认知过程与最终行为应答明确分离。
